\documentclass[12pt,a4paper,oneside]{article}
\usepackage[brazil]{babel}
\usepackage[utf8]{inputenc}
\usepackage{olymp}

\setcounter{problem}{2}

\begin{document}

\begin{problem}{Classificação nas eliminatórias}{classif}{2}
 
A seleção brasileira de futebol tem acumulado uma série de fiascos nas eliminatórias para a próxima Copa do Mundo. Perdeu para Colômbia (pela primeira vez na história das eliminatórias), para o Uruguai (de quem não perdia há 22 anos) e para a Argentina, no Maracanã (pela primeira vez na história das eliminatórias o Brasil perdeu um jogo como mandante). Como se não bastasse, foi eliminada cedo da Copa América deste ano, ainda nas quartas de final. 

Nem a seleção olímpica escapa! Depois do ouro na Rio 2016 e em Tokyo 2020, sequer se classificou para Paris 2024, que começa daqui a algumas horas.

Com os maus resultados nas eliminatórias, está em 6º lugar na América do Sul. Por sorte, ainda na zona de classificação, porque pela primeira vez há 6 vagas, em vez de 4, como de costume.
%se classificarão os melhores 6 times, em vez de 4, como de costume.

Sua tarefa aqui é escrever um programa para ler os resultados das partidas de uma seleção nas eliminatórias da Copa e informar seu total de pontos, seu saldo de gols e seu aproveitamento.

\begin{minipage}{.58\textwidth}
O total de pontos é a soma dos pontos obtidos nas partidas: 3 pts por vitória, 1 pt por empate e 0 por derrota. Por exemplo, nas partidas já disputadas das eliminatórias (resultados ao lado), o Brasil acumulou apenas 7 pontos: 3 + 3 + 1, respectivamente das duas vitórias e um empate nas primeiras partidas.
\end{minipage}
\hfill
\begin{minipage}{.35\textwidth}
\begin{tabular}{rcccl}
\textbf{Brasil} & 5 & x & 1 &Bolívia\\
Peru & 0 & x & 1 & \textbf{Brasil}\\
\textbf{Brasil} & 1 & x & 1 & Venezuela\\
Uruguai & 2 & x & 0 & \textbf{Brasil}\\
Colômbia & 2 & x & 1 & \textbf{Brasil}\\
\textbf{Brasil} & 0 & x & 1 & Argentina
\end{tabular}
\end{minipage}

O saldo de gols é a diferença entre o total de gols marcados e sofridos. Por exemplo, o saldo de gols do Brasil é 1: marcou $5+1+1+0+1+0 = 8$ gols e sofreu $1+0+1+2+2+1 = 7$ gols.

E o aproveitamento é a porcentagem de pontos obtidos do total disputado. Por exemplo, o Brasil disputou 6 partidas. Poderia ter obtido 18 pts mas obteve 7, logo um aproveitamento de 38.9\%.

% Considere uma entrada simplificada, que informa o número de partidas do time e em seguida o resultado de cada partida por dois números, na ordem gols marcados e sofridos pelo time.

%\Notes
%
%Observações, se tiver.

\InputFile

A entrada começa com um número inteiro $N$, o número de partidas disputadas. Em seguida, há $N$ linhas, cada uma com o resultado de uma partida. O resultado de uma partida é dado por uma letra maiúscula e dois inteiros: se a letra for '\texttt{M}', a seleção disputou a partida como mandante, sendo informado primeiro os gols marcados e depois os sofridos; se for '\texttt{V}', disputou como visitante, sendo informado os gols sofridos e depois os marcados. Restrição: $1 \leq N \leq 18$.


\OutputFile

Seu programa deve escrever o total de pontos, o saldo de gols e o aproveitamento da seleção (com 1 casa decimal) nas partidas da entrada, seguindo o modelo dos exemplos a seguir.

% \Notes
% Note que a mensagem escrita não tem acento e tem um ponto final.

\Example

\begin{example}
\exmp{
6
M 5 1
V 0 1
M 1 1
V 2 0
V 2 1
M 0 1
}{
pontos: 7
saldo de gols: 1
aproveitamento: 38.9\%
}%
\end{example}
\textit{\small exemplo do enunciado
}

\begin{example}
\exmp{
6
V 1 0
M 1 0
V 1 1
M 3 0
M 0 0
V 1 1
}{
pontos: 9
saldo de gols: 3
aproveitamento: 50.0\%
}%
\end{example}\\
\textit{\small resultados da Venezuela
}\\

\begin{example}
\exmp{
9
M 5 0
V 2 4
M 1 0
V 0 2
M 2 0
V 0 2
V 0 1
M 2 0
V 1 3
}{
pontos: 27
saldo de gols: 19
aproveitamento: 100.0\%
}%
\end{example}
\textit{\small Brasil nas eliminatóirias da última Copa
}


%Comentários sobre o exemplo, se necessário.

\end{problem}

\end{document}
